% Redefine the abstract environment to align the title and content to the left
\makeatletter
\renewenvironment{abstract}{
    \if@twocolumn
        \section*{\abstractname}%
    \else
        \normalfont\small
        \begin{flushleft} % Change from center to flushleft
            {\bfseries \abstractname\par}%
        \end{flushleft}%
        \noindent\raggedright
    \fi
}
{
    \if@twocolumn\else\par\fi
}
\makeatother

\begin{abstract}

Táto bakalárska práca sa zaoberá spracovaním sekvenačných dát z nanopórového sekvenovania s cieľom 
vytvoriť nástroje a metódy na získanie DNA reťazcov vhodných pre modelovanie narušenia informácie 
v DNA kanáloch. Práca opisuje navrhnutý a implementovaný pipeline, ktorý transformuje surové elektrické 
signály vo formáte pod5 na zarovnané a zklasifikované DNA sekvencie. Kľúčovými komponentmi riešenia 
sú konverzia dát z formátu BAM do textovej podoby, viacnásobné zarovnanie sekvencií pomocou vlastného 
algoritmu inšpirovaného beam search prístupom a identifikácia primeru bez predchádzajúcej znalosti 
jeho presnej sekvencii. Práca demonštruje, že zo surového signálu je možné získať relevantné DNA 
reťazce, ktoré tvoria spoľahlivý základ pre charakterizáciu chybovosti nanopórového čítania. Navrhnutý 
metodický a implementačný rámec umožňuje pokračovať v ďalšom výskume smerom k presnejšiemu modelovaniu 
narušenia informácie v DNA pamätiach.

\smallskip

\noindent\textbf{Kľúčové slová:} nanopórové sekvenovanie, spracovanie DNA dát, zarovnanie sekvencií, DNA pamäť, modelovanie chýb

\end{abstract}
\pagebreak
\begin{otherlanguage}{english}
\begin{abstract}

This bachelor thesis addresses the processing of sequencing data from nanopore sequencing with the goal 
of creating tools and methods to obtain DNA sequences suitable for modeling information perturbation in DNA 
channels. The thesis describes a proposed and implemented pipeline that transforms raw electrical signals 
in pod5 format into aligned and classified DNA sequences. Key components of the solution are conversion of 
data from BAM format to text representation, multiple sequence alignment using a custom algorithm inspired 
by the beam search approach, and primer identification without prior knowledge of its exact sequence. The 
work demonstrates that from raw signal it is possible to obtain relevant DNA sequences that form a reliable 
basis for characterizing the error rate of nanopore reading. The proposed methodological and implementation 
framework enables further research towards more accurate modeling of information perturbation in DNA storage 
systems.

\smallskip

\noindent\textbf{Keywords:} nanopore sequencing, DNA data processing, sequence alignment, DNA storage, error modeling

\end{abstract}
\end{otherlanguage}